\section{Précis}

What follows is a brief précis of the fundamental point of view of the Gornahoor project.

\paragraph{Epistomology}
We hold to three degrees of knowledge:

\begin{itemize}
\item \textbf{Doxa }(\emph{sensus}, opinion, sensual knowledge). Direct, intuitive knowledge through the senses. I taste the sweetness and moistness of a mango. 
\item \textbf{Dianoia }(\emph{ratio}, rational knowledge). Indirect knowledge, discursive reasoning, science. I study the structure of the mango and its history as a crop. 
\item \textbf{Episteme }(\emph{intellectus}, intuition). Direct, intuitive knowledge of the supra-sensible. I grasp the idea of a mango. 
\end{itemize}
\paragraph{Ontology}
\begin{itemize}
\item Spirit is prior to matter 
\item The great chain of being from God, to the spiritual hierarchy, man, animals, plants, matter 
\end{itemize}
\paragraph{Anthropology}
This science describes the nature of man. Man, such as he is, exists in one of three states.

\begin{itemize}
\item The \textbf{sensual man} is focused on \emph{doxa}, or the satisfaction of his material needs and desires. His opinions come ready-made from his social groups and are not questioned or investigated. 
\item The rational man seeks to investigate the sources of the sensual world. There are two options: 

\begin{itemize}
\item \textbf{Science}. Only the sensual world exists. Knowledge is achieved by investigating the facts of the sensual world and forming general laws to explain them. 
\item \textbf{Religion}. A supra-sensible world exists, but there is no way to know it directly. Hence, it is accepted as a matter of belief, based on authority or revelation. 
\end{itemize}
\item The \textbf{intellectual man} has direct access to and knowledge of the supra-sensible realm. 
\end{itemize}
\paragraph{Cosmology}
The sensual world is the reflection of the spiritual world and, thus, has no independent existence. It evolves, in the original sense of the world: as the unfolding of an originary idea. An analogy is that of a symphony, which exists as an idea until musicians perform it; thus the symphony is revealed in the sensual world. As a whole, the world process is the unfolding of the three fundamental ideals. 

\begin{itemize}
\item \textbf{Beauty}. The world process is a creative process, revealing beauty and harmony. 
\item \textbf{Good}. The sufficient reason for the world is the Will. Hence, the world process is a moral action, based on virtue (manliness, strength). The world comes into existence through Power. 
\item \textbf{Truth}. The world is a reflection of the Logos (reason, law). Hence, the world process must follow the logos, or else descend into formlessness. The unity of the logos is highly differentiated and hierarchical, since it is infinite and subject to no limitation. 
\end{itemize}
\paragraph{Politics}
Politics, as the ``queen of the sciences", describes the proper relations of men in their social groupings. As such, it must take into account the findings of all the other sciences, including Anthropology and Cosmology. Hence, all contemporary divisions into left and right are misleading, and ultimately ineffective. 

As Evola pointed out, the correct political position is what ``every well-bred man considered sane, healthy, and normal prior to the French revolution." That is, the harmonious over the discordant, the good over the perverse, order over disorder and the undifferentiated mass. 

The battle, therefore, is between those who would uphold the Traditional order, and the various revolutionary forces that seek to overthrow it. The revolution operates on multiple fronts. 

\paragraph{Western Tradition}
The Western Tradition is based on the axis from Greece and Rome, and its prolongation in time to Paris (Charles Maurras). We would extend that back even further to its Indo-European roots in Egypt, Persia, India. We, therefore, see continuity from the classical period to the Catholic medieval era and even into modern times. 

We acknowledge the existence of a Western esoteric tradition, which at various times has made itself known and at others, has had to hide itself. This tradition has been based on Hermeticism, whether the ancient pagan mysteries or Christian alchemy. The goal of this Tradition is the reintegration of man into his Primordial state; this is true, whether pagan or Christian. 

Hence, we reject the typical superficial distinctions between pagan and Christian. This requires the simultaneous rejection of those vulgar expressions of paganism as the life of non-transcendent sensuality, and Christian other-worldliness unnourished by its esoteric sources.



\flrightit{Posted on 2010-05-08 by Cologero }

\begin{center}* * *\end{center}

\begin{footnotesize}\begin{sffamily}



\texttt{Matt on 2010-05-08 at 21:15 said: }

In thinking about the Western Tradition, and the Tradition in general, I have issues about where Sufism falls. The fundamental point in the Tradition (correct me if I'm wrong) is that the pure Self (not the everyday ego) must be realized and affirmed. And this fundamental point is stated in both the West (as seen in Hermeticism and the mystery schools), and the East (the Buddha's original doctrine, Vedanta and Yoga, Taosim) with the difference being how the sensible world should be treated. Yet with the Sufi view, it appears that the Self must not be affirmed, but annihilated in the divine (granted, Sufism may be alluding to the physical ego, and with going forward in studying it, that could very well be the case). So at least right now, I am perplexed as to where this Islamic mysticism falls under in the Tradition.


\hfill

\texttt{Francis Mercuri on 2010-05-09 at 02:11 said: }

This ``Precis" of the Gornahoor project is an inspirational and succinct statement of principles; your work and writings are appreciated.

Additionally, for those following a Christian path, this tertiary anthropology might also be extended to the correspondences with the ``degrees" of the primitive Christian hierarchy:

1. Catechumens

2. The Faithful

3. Christians 

Somewhat analogous to the Hellenistic:

1. Hylics

2. Psychics

3. Pneumatics

The Masonic:

1. Apprentices

2. Fellowcrafts

3. Masters

Mouravieff's:

1. Exoteric cycle

2. Mesoteric cycle

3. Esoteric cycle

And Robin Amis’ ``Three Renunciations", of:

1. Physical ascesis

2. Psychological ascesis

3. Noetic ascesis, which his exegesis compares with:

1. Proverbs

2. Ecclesiastes

3. The Song of Songs


\hfill

\texttt{Francis Mercuri on 2010-05-10 at 04:43 said: }

Mark:

1. I'm in agreement with some others who say that Islam (and thus Sufism), is neither really ``Eastern" nor ``Western"; instead, the Islamic tradition represents something of a ``amalgam" path, or a liaison between East and West.

2. Along the lines of Guenon's definitions, I don't perceive Sufism as a type of ``mysticism", as you call it. Guenon lucidly explained the differences between Initiation and mysticism in his ``Perspectives on Initiation", where, in brief, he noted that mysticism always remains an aspect of exotericism (and the theological domain), in that in the ``mystical" experience, the individual still remains the same individual, albeit in a state of passive reception of Superior influences, which are always ``colored" by a given theological clothing, and occur in a rather random way. 

On the other hand, the Initiate goes beyond the theological domain, meeting the metaphysical, and intentionally seeks simulatenously to divinize his individuality, and individualize the Personality…this is what the Sufis call the ``Supreme Identity", attained in a conscious and permanent manner. While anybody can be a ``mystic", the Initiate is set apart by specific ``qualifications", and travels a well governed path, which is first defined by ``psychic regeneration" (the ``Primordial State"), and then, if possible Noetic regeneration (the ``Transcendent State"). Echoes of this can be noted in the Western tradition, where a man being made a Priest by a Bishop is said to be ``elevated" to the Priesthood, or a Mason, on becoming a Master is said to be ``raised" to the Sublime Degree . The forum owner's recent explanations of the three types of epistemology, are another way of expressing the same things. 

3. In Sufism, there is indeed, as you observe a emphasis on ``extinction" (Fana), and beyond that ``extinction of extinction" (Fana al-Fana). But, I'm not so sure that I'd agree that this implies a destruction of the individuality, so much as an ``extinction" of what the Sufi psychology calls the Nafs (which are what we previously identified as the Christian ``prilog", ``passions", or in any other words, false ``i's"). After all, the state of ``Fana" corresponds to the grade of Insan al-Kadim (Primordial man), while the higher (highest) of Fana al-Fana corresponds with Insan al-Kamil (Universal/Transcendent Man). Were this not so, the Prophet Himself, who is regarded as an exemplar of Insan al-Kamil, could not be spoken of as having any individuality! No, what is involved here, as explained in section 2, is an example of how traditional Initiation operates:

a. The Prophet as Insan al-Kamil, is equivalent to:

b. The ``Nur" Mohammed

c. The ``Transcendent Man"

d. The regenerated Nous, or Logos

e. The esoteric Buddhist ``Dianchi"

f. The ``Celestial Jerusalem"

g. The epistomological ``Episteme"

h. The ``Third Birth", aka ``Resurrection"

i. The Hermetic ``projection"

j. The Hindu ``Avatar'

k. The Masonic ``Master"

l. The Tibetan ``Lama"

m. etc……


\hfill

\texttt{Matt on 2010-05-10 at 12:01 said: }

I appreciated that response Francis.

Yes, I probably should have used initiatic teaching instead of mysticism. Some of what I have read about Sufism has probably not been from the best sources and as stated before, I will have to further study it. What you stated about fana has helped my insight though. Thank you.

Oh,and you got me confused with Mark. I'm Matt.


\hfill

\texttt{Roger Buck on 2010-05-17 at 15:27 said: }

Very good to have this! For now I will only comment on one most striking phrase: ``its prolongation in time to Paris"!

I don't really know Maurras, but find myself wanting to know more of either why he says this and/or why you Cologero include this in this précis.

It is quite stirring and moving for me, this phrase and I cannot help but wonder what connexions it might possibly have with the reasons why Meditations on the Tarot was written in French, and sought to incarnate into a stream of French Hermeticists who at least in some cases, would have shared similar notions.

Yes I wonder why is Paris singled out like this in the Gornahoor précis …


\hfill

\texttt{Cologero on 2010-05-22 at 11:42 said: }

We accept the general approach of Marcilio Ficino: Hermes $\rightarrow$ Pythagoras $\rightarrow$ Plato $\rightarrow$ Neoplatonism (simplified). After Ficino, we find the fulfillment of the Hermetic Tradition among the French: Louis-Claude de St-Martin, Joseph de Maistre, up to the French esoterists of the late 19th and early 20th centuries.


\end{sffamily}\end{footnotesize}
